\documentclass[10pt]{article}
\usepackage[T1]{fontenc}
\usepackage[utf8]{inputenc}
\usepackage[french]{babel}
\usepackage{amsmath,amssymb,mathrsfs,mathpazo}
\usepackage{geometry}
\geometry{a4paper,margin=2cm}
\usepackage{graphicx,array,tabularx,tabularray,colortbl,multirow,multicol,booktabs,enumitem,pifont,ulem,fancybox,diagbox,cancel}
\usepackage{pstricks,pstricks-add,pst-plot,pst-tree,pst-node,pst-text,pst-func,pst-math,pst-bspline,pst-3dplot,pst-eucl,pst-solides3d}
\let\pscalculate\undefined
\usepackage{pst-bezier}
\usepackage{tikz,pgfplots,tkz-tab}
\usetikzlibrary{arrows,arrows.meta,patterns,calc,positioning,trees,decorations.markings,shapes.geometric,shapes.misc,shapes.arrows,intersections,decorations.pathmorphing,decorations.pathreplacing,backgrounds,shadows,babel}
\pgfplotsset{compat=1.18}
\usepgfplotslibrary{fillbetween}
\usepackage[np]{numprint}
\usepackage{listings,makecell,pas-tableur,twoopt}
\let\barre\undefined
\usepackage[pstricks]{tablvar}
\definecolor{darkgreen}{RGB}{0,100,0}
\providecommand{\red}{\color{red}}
\providecommand{\blue}{\color{blue}}
\providecommand{\green}{\color{green}}
\providecommand{\black}{\color{black}}
\providecommand{\white}{\color{white}}
\providecommand{\gray}{\color{gray}}
\providecommand{\cyan}{\color{cyan}}
\providecommand{\magenta}{\color{magenta}}
\providecommand{\yellow}{\color{yellow}}
\providecommand{\violet}{\color{violet}}
\providecommand{\orange}{\color{orange}}
\providecommand{\euro}{\texteuro}
\providecommand{\EUR}{\texteuro}
\providecommand{\ieme}{\textsuperscript{e}}
\providecommand{\degres}{\ensuremath{^\circ}}
\providecommand{\degree}{\ensuremath{^\circ}}
\providecommand{\vv}[1]{\overrightarrow{#1}}
\providecommand{\Coord}[1]{#1}
\providecommand{\decosix}{}
\providecommand{\starredbullet}{\ensuremath{\bullet}}
\pagestyle{empty}
\setlength{\parindent}{0pt}
\usepackage{adjustbox}\DeclareUnicodeCharacter{2236}{:}\DeclareUnicodeCharacter{2190}{\ensuremath{\leftarrow}}
\providecommand{\R}{\mathbb{R}}
\providecommand{\N}{\mathbb{N}}
\providecommand{\Z}{\mathbb{Z}}
\providecommand{\Q}{\mathbb{Q}}
\providecommand{\D}{\mathbb{D}}
\providecommand{\C}{\mathbb{C}}
\providecommand{\e}{{\mathrm{e}}}
\providecommand{\vect}[1]{\overrightarrow{#1}}
\providecommand{\vectt}[1]{\overrightarrow{\mathrm{#1}}}
\providecommand{\vvt}[1]{\overrightarrow{\mathrm{#1}}}
\providecommand{\np}[2][]{{#2~#1}}
\providecommand{\vv}[1]{{\overrightarrow{#1}}}
\renewcommand{\d}{{\mathrm{d}}}
\providecommand{\degres}{{^\circ}}
\providecommand{\degre}{{^\circ}}
\providecommand{\blue}{\color{blue}}
\providecommand{\red}{\color{red}}
\providecommand{\enskip}{\quad}
\providecommand{\llbracket}{\lbrack\!\lbrack}
\providecommand{\rrbracket}{\rbrack\!\rbrack}
\providecommand{\euro}{{\text{€}}}
\providecommand{\no}{{\text{n°}}}
\providecommand{\interval}[2]{{[#1;#2]}}
\providecommand{\intervFO}[2]{{[#1;#2[}}
\providecommand{\dotfill}{\ldots}

\providecommand{\ds}{\displaystyle}
\providecommand{\textding}[1]{#1}
\providecommand{\Oij}{\(\mathrm{O};\vec{\imath},\vec{\jmath}\)}
\providecommand{\Oijk}{\(\mathrm{O};\vec{\imath},\vec{\jmath},\vec{k}\)}
\providecommand{\Ouv}{\(\mathrm{O};\vec{u},\vec{v}\)}
\providecommand{\intervFF}[2]{\ensuremath{[#1;#2]}}
\renewcommand{\euro}{\ifmmode\text{\texteuro}\else\texteuro\fi}
\DeclareUnicodeCharacter{00B0}{\ifmmode^{\circ}\else\textdegree\fi}
\DeclareUnicodeCharacter{20AC}{\ifmmode\text{\texteuro}\else\texteuro\fi}
% Énoncés : Éducation nationale. Transcription et corrigés : APMEP, Denis Vergès et contributeurs. Découpage et classement : CoopMaths, Sébastien Lozano, Éric Elter et contributeurs. Adaptation de la présentation : Maths Explorer.
% CC BY-SA : https://coopmaths.fr/www/about/
% https://coopmaths.fr/alea/static/eam/2026/tex/eam_specialite_2026_06_polynesie_automatismes.tex
\begin{document}
\noindent\textit{APMEP / CoopMaths. CC BY-SA. Présentation adaptée par Maths Explorer.}\par\medskip
\textbf{Pour cette première partie, aucune justification n'est demandée et une seule réponse est possible par question.\\
	Pour chaque question, reportez son numéro sur votre copie et indiquez votre réponse.}

	Une réponse fausse ou l'absence de réponse n'enlève aucun point.

\subsubsection*{Question 1}
Un article coûte 40~\euro{} et subit une augmentation de 30\,\%.

Combien coûte-t-il après cette augmentation ?
		\begin{enumerate}[label=\alph*)]
\begin{multicols}{4}
			\item \np{40.3}~\euro
			\item \np{52}~\euro
			\item \np{12}~\euro
			\item \np{70}~\euro
				\end{multicols}
		\end{enumerate}

		\subsubsection*{Question 2}

	On considère l'équation \mbox{$(-\np{0.5}x+3)(-5x-4)=0$}.

	Les solutions de cette équation sont :
		\begin{enumerate}[label=\alph*)]
			\begin{multicols}{4}
				\item \mbox{$\dfrac{\np{0.5}}{3}\text{ et }\dfrac54$}
				\item \mbox{$\dfrac{\np{0.5}}{3}\text{ et }-\np{0.8}$}
				\item \mbox{$6\text{ et }\np{0.8}$}
				\item \mbox{$6\text{ et }-\np{0.8}$}

				\end{multicols}
		\end{enumerate}

\subsubsection*{Question 3}
Un entraîneur choisit 8 joueurs dans une équipe.

Cela représente 20\,\% de l'équipe.

Combien y a-t-il de joueurs dans l'équipe ?
\begin{enumerate}[label=\alph*)]
	\begin{multicols}{4}
		\item \mbox{$16$}
		\item \mbox{$20$}
		\item \mbox{$40$}
		\item \mbox{$32$}

		\end{multicols}
\end{enumerate}



\subsubsection*{Question 4}
On peut calculer l'énergie cinétique d'un objet en mouvement.

Cette énergie est notée \mbox{$E_C$} et elle est donnée par la formule \mbox{$E_C=\dfrac12~m~v^2$} (où $E_C$ s'exprime en joules, $m$ représente la masse de l'objet en kg et $v$ sa vitesse en m/s).

\vspace{5mm}
L'expression permettant d'exprimer la vitesse $v$ en fonction de $E_C$ et de $m$ est :


\begin{enumerate}[label=\alph*)]
	\begin{multicols}{4}
		\item \mbox{$v=\dfrac{{E_C} ^2}{2m}$}
		\item \mbox{$v=\sqrt{\dfrac{2~E_C }{m}}$}
		\item \mbox{$v=\sqrt{\dfrac{E_C }{2~m}}$}
		\item \mbox{$v=\sqrt{2~m~E_C }$}
				\end{multicols}
\end{enumerate}
\subsubsection*{Question 5}
\begin{minipage}[t]{0.5\linewidth}\vspace{0mm}
On a représenté ci-contre la courbe de la fonction inverse définie sur $\mathbb{R}$* par \mbox{$f(x)=\dfrac1 x$.}


L'ensemble des solutions de l'inéquation \mbox{$\dfrac1 x \geqslant  2$} est :
	\begin{enumerate}[label=\alph*)]
		\begin{multicols}{2}
			\item \mbox{$\left]0;\dfrac12\right]$}
			\item \mbox{$\left]0;2\right]$}

		\end{multicols}
		\begin{multicols}{2}
			\item \mbox{$\left[2;+\infty\right[$}
			\item \mbox{$\left]-\infty;0 \vphantom{\dfrac{1}{2}}\right[\cup\left[\dfrac12;+\infty\right[$}
		\end{multicols}
	\end{enumerate}
\end{minipage}\hfill
	\begin{tikzpicture}[baseline={(current bounding box.north)},xscale=1,yscale=1]
		\clip(-3.4,-3.4)rectangle(3.4,3.4);
		\draw[blue!20, very thin](-3.4,-3.4)grid[xstep=0.2,ystep=0.2](3.4,3.4);
		\draw[blue!30,thick](-3.4,-3.4)grid(3.4,3.4);
		\draw[black, line width=0.3mm,->, >=triangle 45](-3.4,0)--(3.4,0);
		\draw[black, line width=0.3mm,->, >=triangle 45](0,-3.4)--(0,3.4);
	\foreach \i in {-3,-2,-1,1,2,3}
	{\draw (\i,2pt)--(\i,-2pt)node[below]{\np{\i}} (2pt,\i)(-2pt,\i)node[left]{\np{\i}};}
	\draw(-0.2,-0.2)node{\np{0}};
	\draw[red, line width=0.3mm]plot[smooth,samples=50, domain=-3.4:-0.1](\x,{1/\x});
	\draw[red, line width=0.3mm]plot[smooth,samples=50, domain=0.1:3.4](\x,{1/\x});
		\end{tikzpicture}

\subsubsection*{Question 6}
Le nombre \mbox{$2^9\times5^7$} est égal à :
		\begin{enumerate}[label=\alph*)]
	\begin{multicols}{4}
		\item \mbox{$10^{16}$}
		\item \mbox{$4\times10^7$}
		\item \mbox{$10^{63}$}
		\item \mbox{$4\times10^{14}$}
		\end{multicols}
\end{enumerate}
\subsubsection*{Question 7}
\begin{minipage}[t]{0.45\linewidth}\vspace{0mm}
	Parmi les équations réduites de droites proposées, laquelle est celle de la droite $(d)$ tracée dans le repère ci-contre?

		\begin{enumerate}[label=\alph*)]
		\begin{multicols}{2}
			\item \mbox{$y=-\np{5.2}x+\np{8.6}$}
			\item \mbox{$y=\np{0.6}x+\np{5.2}$}

		\end{multicols}
		\begin{multicols}{2}
			\item \mbox{$y=\np{5.2}x+\np{8.6}$}
			\item \mbox{$y=-\np{0.6}x+\np{5.2}$}

		\end{multicols}
	\end{enumerate}

\end{minipage}\hfill
	\begin{tikzpicture}[baseline={(current bounding box.north)},xscale=0.6,yscale=0.6,rotate=0]
		\clip(-2.3,-.65)rectangle(9.2,6.4);
		\draw[blue!30,thick](-2.3,-0.5)grid(9.2,6.4);
		\draw[line width=0.3mm,->, >=triangle 45](-2.3,0)--(9.2,0);
		\draw[line width=0.3mm,->, >=triangle 45](0,-0.5)--(0,6.4);
		\foreach \i in {-3,-2,-1,1,2,3,4,5,6,7,8,9}
		{\draw (\i,2pt)(\i,-2pt)node[below]{\np{\i}} (2pt,\i)--(-2pt,\i)node [left]{\np{\i}};}
		\draw(-7pt,-2 pt)node[below]{\np{0}};
		\draw[red, line width=0.3mm](-2.3,6.58)--(9.2,-0.32)node[pos=0.1,fill=white, below left]{\textbf{$(d)$}};
	\end{tikzpicture}

\subsubsection*{Question 8}
Quelle est la forme factorisée de l'expression \mbox{$16x^2-(x+1)^2$} ?

\begin{enumerate}[label=\alph*)]
	\begin{multicols}{4}
		\item \mbox{$(3x-1)^2$}
		\item \mbox{$(3x-1)(5x+1)$}
		\item \mbox{$(15x-1)^2$}
		\item \mbox{$(15x-1)(17x+1)$}
		\end{multicols}
\end{enumerate}

\end{document}
